% !PRINTED 627
\weberpage{627}
Der Übergang von einer Basis $\eta_1, \eta_2, \ldots, \eta_n$ zu einer anderen $\eta_1', \eta_2', \ldots, \eta_n'$ geschieht durch eine lineare Substitution:
\weq{9}{(\eta_1', \eta_2', \ldots, \eta_n') = L(\eta_1, \eta_2, \ldots, \eta_n)}
oder kürzer geschrieben
\weq{10}{(\eta_i') = L(\eta_i),}
worin $L$ die Bedeutung hat
\weq{11}{
L = \begin{pmatrix}
l_{1,1}, & l_{1,2}, & \ldots, & l_{1,n} \\
l_{2,1}, & l_{2,2}, & \ldots, & l_{2,n} \\
\cdot & \cdot & \cdot & \cdot \\
l_{n,1}, & l_{n,2}, & \ldots, & l_{n,n}
\end{pmatrix} \quad \text{(Bd.\,II, \S\,41).}
}

Die Elemente $l_{i,k}$ dieser Substitution sind rationale Funktionen von $z$.

\gesperrt{Das System $(\eta')$ ist dann und nur dann wieder eine Basis, wenn die Determinante $|L|$ dieser Substitution nicht verschwindet.}

\begin{center}
\S\ 172.

\gesperrt{Normen und Spuren.}
\end{center}

Wählt man zur Darstellung der Funktionen in $\Omega$ eine beliebige Basis $\eta_1, \eta_2, \ldots, \eta_n$, so kann man, wenn $\zeta$ eine beliebige Funktion des Körpers ist, da die Produkte $\zeta\eta_i$ im Körper enthalten sind, setzen:
\weq{1}{
\begin{aligned}
\zeta\eta_1 &= y_{1,1}\eta_1 + y_{1,2}\eta_2 + \cdots + y_{1,n}\eta_n, \\
\zeta\eta_2 &= y_{2,1}\eta_1 + y_{2,2}\eta_2 + \cdots + y_{2,n}\eta_n, \\
&\qquad\qquad\qquad\cdots\cdots\cdots\cdots \\
\zeta\eta_n &= y_{n,1}\eta_1 + y_{n,2}\eta_2 + \cdots + y_{n,n}\eta_n,
\end{aligned}
}
und daraus erhält man durch Elimination der $\eta$:
\weq{2}{
\begin{vmatrix}
y_{1,1} - \zeta, & y_{1,2}, & \ldots, & y_{1,n} \\
y_{2,1}, & y_{2,2} - \zeta, & \ldots, & y_{2,n} \\
\cdot & \cdot & \cdot & \cdot \\
y_{n,1}, & y_{n,2}, & \ldots, & y_{n,n} - \zeta
\end{vmatrix} = 0,
}
eine Gleichung, die nach Potenzen von $\zeta$ geordnet, die Gestalt erhält:
\weq{3}{\varphi(\zeta) = \zeta^n + b_1\zeta^{n-1} + \cdots + b_{n-1}\zeta + b_n = 0,}
in der $b_1, b_2, \ldots, b_n$ rationale Funktionen von $z$ sind. Diese Koeffizienten sind von der Wahl der Basis $\eta_i$ unabhängig. Denn die Gleichungen (1), durch die der Übergang von der Basis $\eta_i$
%CONT>>
